\documentclass[a4paper,twoside]{article}

\usepackage{morewrites}

\usepackage[svgnames,dvipsnames,x11names]{xcolor}
\usepackage{graphicx}
\usepackage[english]{babel}
\usepackage[colorlinks=true, linkcolor=NavyBlue, urlcolor=NavyBlue, citecolor=NavyBlue]{hyperref}
\usepackage[a4paper]{geometry}
\usepackage{ts-fonts}
\usepackage{amsmath}
\usepackage{float}
\usepackage{seqsplit}
\usepackage{ts-callouts}
\usepackage{ts-code}

\usepackage{lastpage}
\usepackage{refcount}
\makeatletter
\def\ps@plain{
  \let\@oddhead\@empty
  \let\@evenhead\@empty
  \def\@oddfoot{
    \hfil
    \ifnum\getpagerefnumber{LastPage}>1\relax
      \thepage
    \fi
    \hfil}
  \let\@evenfoot\@oddfoot
}
\makeatother

\title{YAML Front Matter}
\author{}\date{}

\begin{document}

\maketitle

\thispagestyle{plain}
Every Markdown document can carry a small block of metadata at the very top, fenced by \texttt{-\allowbreak{}-\allowbreak{}-\allowbreak{}} markers. This is the \textbf{front matter}: a YAML island living rent-free above your prose. Static site generators like MkDocs read it to drive per-page configuration, and TeXSmith hooks into the same convention to steer how a document is parsed, typeset, and ultimately rendered to \LaTeX{} or PDF.

A guiding principle: \textbf{keep content and form apart}. Templates, font sizes, margins, paper format, and other typographic knobs belong in the front matter; the body should care only about ideas, sentences, and equations. Other tools take a different stance, see \href{https://quarkdown.com/}{Quarkdown}, which weaves configuration directly into the document body. Both are valid, but TeXSmith favors the separation, your future self will thank you when swapping templates without touching a single paragraph.

\section{Press}\label{press}

The \texttt{press} block holds everything related to the printed (or PDF'd) artifact, title, authors, template choice, and any template-specific slots:

\begin{code}{yaml}{}{}
\begin{Verbatim}[commandchars=\\\{\},breaklines, breakanywhere, commandchars=\\\{\}]
\PY{n+nt}{press}\PY{p}{:}
\PY{+w}{  }\PY{n+nt}{title}\PY{p}{:}\PY{+w}{ }\PY{l+s}{\PYZdq{}}\PY{l+s}{My}\PY{n+nv}{ }\PY{l+s}{Document}\PY{n+nv}{ }\PY{l+s}{Title}\PY{l+s}{\PYZdq{}}
\PY{+w}{  }\PY{n+nt}{subtitle}\PY{p}{:}\PY{+w}{ }\PY{l+s}{\PYZdq{}}\PY{l+s}{An}\PY{n+nv}{ }\PY{l+s}{In\PYZhy{}depth}\PY{n+nv}{ }\PY{l+s}{Exploration}\PY{l+s}{\PYZdq{}}
\PY{+w}{  }\PY{n+nt}{template}\PY{p}{:}\PY{+w}{ }\PY{l+lScalar+lScalarPlain}{article}
\PY{+w}{  }\PY{n+nt}{authors}\PY{p}{:}
\PY{+w}{    }\PY{p+pIndicator}{\PYZhy{}}\PY{+w}{ }\PY{n+nt}{name}\PY{p}{:}\PY{+w}{ }\PY{l+s}{\PYZdq{}}\PY{l+s}{Alice}\PY{n+nv}{ }\PY{l+s}{Smith}\PY{l+s}{\PYZdq{}}
\PY{+w}{      }\PY{n+nt}{affiliation}\PY{p}{:}\PY{+w}{ }\PY{l+s}{\PYZdq{}}\PY{l+s}{University}\PY{n+nv}{ }\PY{l+s}{of}\PY{n+nv}{ }\PY{l+s}{Examples}\PY{l+s}{\PYZdq{}}
\PY{+w}{  }\PY{n+nt}{slots}\PY{p}{:}
\PY{+w}{    }\PY{n+nt}{abstract}\PY{p}{:}\PY{+w}{ }\PY{l+lScalar+lScalarPlain}{Abstract}
\end{Verbatim}
\end{code}
\begin{callout}[callout note]{Note}
The \texttt{press} section is optional. Keys like \texttt{title} and \texttt{authors} may also live at the root of the front matter, and TeXSmith will pick them up just fine. Nesting them under \texttt{press} keeps things tidy and avoids stepping on the toes of other static site generators that may want the root-level keys for themselves.
\end{callout}
Each template exposes its own set of attributes (cover styles, sidebar toggles, custom slots, ...). Head over to the \hyperref[snippet:<HASH>]{Template Guide} for the full menu.

\section{Bibliography}\label{bibliography}

References can be declared inline, right next to the document that cites them, no external \texttt{.bib} file required (though one still works if you prefer). Mix DOI shortcuts with fully-spelled-out entries as needed:

\begin{code}{yaml}{}{}
\begin{Verbatim}[commandchars=\\\{\},breaklines, breakanywhere, commandchars=\\\{\}]
\PY{n+nt}{bibliography}\PY{p}{:}
\PY{+w}{  }\PY{n+nt}{AB2020}\PY{p}{:}\PY{+w}{ }\PY{l+lScalar+lScalarPlain}{doi:10.1000/xyz123}
\PY{+w}{  }\PY{n+nt}{CD2019}\PY{p}{:}
\PY{+w}{    }\PY{n+nt}{type}\PY{p}{:}\PY{+w}{ }\PY{l+lScalar+lScalarPlain}{book}
\PY{+w}{    }\PY{n+nt}{author}\PY{p}{:}\PY{+w}{ }\PY{l+s}{\PYZdq{}}\PY{l+s}{John}\PY{n+nv}{ }\PY{l+s}{Doe}\PY{l+s}{\PYZdq{}}
\PY{+w}{    }\PY{n+nt}{title}\PY{p}{:}\PY{+w}{ }\PY{l+s}{\PYZdq{}}\PY{l+s}{Example}\PY{n+nv}{ }\PY{l+s}{Book}\PY{l+s}{\PYZdq{}}
\PY{+w}{    }\PY{n+nt}{year}\PY{p}{:}\PY{+w}{ }\PY{l+s}{\PYZdq{}}\PY{l+s}{2019}\PY{l+s}{\PYZdq{}}
\end{Verbatim}
\end{code}
The full syntax, supported entry types, and resolution rules are documented in the \hyperref[snippet:<HASH>]{Bibliography Guide}.

\section{Glossary}\label{glossary}

When the glossary feature is enabled, entries are declared in the front matter and grouped into logical tables. Symbols, acronyms, and domain-specific jargon all coexist peacefully:

\begin{code}{yaml}{}{}
\begin{Verbatim}[commandchars=\\\{\},breaklines, breakanywhere, commandchars=\\\{\}]
\PY{n+nt}{glossary}\PY{p}{:}
\PY{+w}{  }\PY{n+nt}{style}\PY{p}{:}\PY{+w}{ }\PY{l+lScalar+lScalarPlain}{long}\PY{+w}{ }\PY{c+c1}{\PYZsh{} or short}
\PY{+w}{  }\PY{n+nt}{groups}\PY{p}{:}\PY{+w}{ }\PY{c+c1}{\PYZsh{} Grouping in different tables (optional)}
\PY{+w}{    }\PY{n+nt}{symbols}\PY{p}{:}\PY{+w}{ }\PY{l+lScalar+lScalarPlain}{Mathematical}\PY{l+lScalar+lScalarPlain}{ }\PY{l+lScalar+lScalarPlain}{symbols}\PY{l+lScalar+lScalarPlain}{ }\PY{l+lScalar+lScalarPlain}{and}\PY{l+lScalar+lScalarPlain}{ }\PY{l+lScalar+lScalarPlain}{notations}
\PY{+w}{    }\PY{n+nt}{corporate}\PY{p}{:}\PY{+w}{ }\PY{l+lScalar+lScalarPlain}{Organizational}\PY{l+lScalar+lScalarPlain}{ }\PY{l+lScalar+lScalarPlain}{terms}
\PY{+w}{    }\PY{n+nt}{technology}\PY{p}{:}\PY{+w}{ }\PY{l+lScalar+lScalarPlain}{Technology\PYZhy{}related}\PY{l+lScalar+lScalarPlain}{ }\PY{l+lScalar+lScalarPlain}{terms}
\PY{+w}{  }\PY{n+nt}{entries}\PY{p}{:}
\PY{+w}{    }\PY{l+s}{\PYZdq{}}\PY{l+s}{\PYZdl{}}\PY{l+s}{\PYZbs{}\PYZbs{}}\PY{l+s}{phi\PYZdl{}}\PY{l+s}{\PYZdq{}}\PY{p+pIndicator}{:}
\PY{+w}{      }\PY{n+nt}{group}\PY{p}{:}\PY{+w}{ }\PY{l+lScalar+lScalarPlain}{symbols}
\PY{+w}{      }\PY{n+nt}{description}\PY{p}{:}\PY{+w}{ }\PY{l+lScalar+lScalarPlain}{Angle}\PY{l+lScalar+lScalarPlain}{ }\PY{l+lScalar+lScalarPlain}{in}\PY{l+lScalar+lScalarPlain}{ }\PY{l+lScalar+lScalarPlain}{radians}
\PY{+w}{    }\PY{n+nt}{ONU}\PY{p}{:}
\PY{+w}{      }\PY{n+nt}{group}\PY{p}{:}\PY{+w}{ }\PY{l+lScalar+lScalarPlain}{corporate}
\PY{+w}{      }\PY{n+nt}{description}\PY{p}{:}\PY{+w}{ }\PY{l+lScalar+lScalarPlain}{United}\PY{l+lScalar+lScalarPlain}{ }\PY{l+lScalar+lScalarPlain}{Nations}\PY{l+lScalar+lScalarPlain}{ }\PY{l+lScalar+lScalarPlain}{Organization}
\PY{+w}{    }\PY{n+nt}{AI}\PY{p}{:}
\PY{+w}{      }\PY{n+nt}{group}\PY{p}{:}\PY{+w}{ }\PY{l+lScalar+lScalarPlain}{technology}
\PY{+w}{      }\PY{n+nt}{description}\PY{p}{:}\PY{+w}{ }\PY{l+lScalar+lScalarPlain}{Artificial}\PY{l+lScalar+lScalarPlain}{ }\PY{l+lScalar+lScalarPlain}{Intelligence}
\end{Verbatim}
\end{code}
See the \hyperref[snippet:<HASH>]{Glossary Guide} for sorting behavior, cross-references, and styling options.

\section{Counters}\label{counters}

Document-specific numbered series --- findings, requirements, bugs, test cases ---
are declared under \texttt{counters:}. Each key is the prefix used by the \texttt{\#\{…\}}
markers and the \texttt{@…} references in the body:

\begin{code}{yaml}{}{}
\begin{Verbatim}[commandchars=\\\{\},breaklines, breakanywhere, commandchars=\\\{\}]
\PY{n+nt}{counters}\PY{p}{:}
\PY{+w}{  }\PY{n+nt}{n}\PY{p}{:}
\PY{+w}{    }\PY{n+nt}{name}\PY{p}{:}\PY{+w}{ }\PY{l+lScalar+lScalarPlain}{Requirement}\PY{+w}{ }\PY{c+c1}{\PYZsh{} human\PYZhy{}readable name, used in diagnostics}
\PY{+w}{    }\PY{n+nt}{format}\PY{p}{:}\PY{+w}{ }\PY{l+s}{\PYZdq{}}\PY{l+s}{N\PYZhy{}\PYZob{}n:02d\PYZcb{}}\PY{l+s}{\PYZdq{}}\PY{+w}{ }\PY{c+c1}{\PYZsh{} optional, defaults to \PYZdq{}\PYZob{}n\PYZcb{}\PYZdq{}}
\PY{+w}{    }\PY{n+nt}{start}\PY{p}{:}\PY{+w}{ }\PY{l+lScalar+lScalarPlain}{1}\PY{+w}{ }\PY{c+c1}{\PYZsh{} optional, defaults to 1}
\PY{+w}{  }\PY{n+nt}{fw}\PY{p}{:}
\PY{+w}{    }\PY{n+nt}{name}\PY{p}{:}\PY{+w}{ }\PY{l+lScalar+lScalarPlain}{Firmware}\PY{l+lScalar+lScalarPlain}{ }\PY{l+lScalar+lScalarPlain}{finding}
\PY{+w}{    }\PY{n+nt}{format}\PY{p}{:}\PY{+w}{ }\PY{l+s}{\PYZdq{}}\PY{l+s}{FW\PYZhy{}\PYZob{}n:02d\PYZcb{}}\PY{l+s}{\PYZdq{}}
\end{Verbatim}
\end{code}
See \hyperref[snippet:<HASH>]{Custom counters} for the marker syntax, the
numbering scope and the backend mapping.

\section{Cross-document references}\label{cross-document-references}

\texttt{id} gives the document a free identifier (a contract or report number);
\texttt{crossrefs:} declares the inventories published by the documents it cites:

\begin{code}{yaml}{}{}
\begin{Verbatim}[commandchars=\\\{\},breaklines, breakanywhere, commandchars=\\\{\}]
\PY{n+nt}{id}\PY{p}{:}\PY{+w}{ }\PY{l+lScalar+lScalarPlain}{RHE\PYZhy{}424}
\PY{n+nt}{crossrefs}\PY{p}{:}
\PY{+w}{  }\PY{n+nt}{fwrev}\PY{p}{:}\PY{+w}{ }\PY{l+lScalar+lScalarPlain}{build/firmware\PYZhy{}review.refs.json}
\end{Verbatim}
\end{code}
A citation then reads \texttt{@fwrev:fw:pas-\allowbreak{}de-\allowbreak{}temps} and renders as
\texttt{RHE-\allowbreak{}423-\allowbreak{}FW-\allowbreak{}10 p. 14}. See \hyperref[snippet:<HASH>]{Cross-document references}.

\end{document}
